\documentclass[10pt,a4paper]{amsart}
\usepackage[margin=19mm]{geometry}
\usepackage{amsmath,amssymb,mathtools}
\usepackage[scr=rsfs,cal=euler]{mathalfa}
\usepackage{xfrac}
\usepackage[unicode,hidelinks,bookmarks=false]{hyperref}
\newcommand{\ie}{i.e.\ }
\newcommand{\cf}{cf.\ }
\newcommand{\resp}{respectively\ }
\newcommand{\Subsection}{\subsection}
\providecommand{\nobreakdash}{\nobreak\hbox{-}}
\setlength{\emergencystretch}{2em}
\multlinegap0pt
\usepackage{amssymb}
\usepackage{mathtools}
\usepackage{tikz}
\newtheorem{lemma}{Lemma}
\newcommand{\CC}{{\mathbb C}}
\newcommand\R{\mathbb R}
\newcommand{\RR}{{\mathbb R}}
\newcommand\bX{{\mathbb X}}
\newcommand\cT{{\mathcal T}}
\newcommand{\dom}{\text{\rm{dom}}}
\newcommand{\rmd}{{\mathrm{d}}}
\title{The Damped Waves Equation and generalized Cosine and Sine families on Banach spaces}
\author{Yuri Latushkin, Alin Pogan}
\date{Source: arXiv:2607.05856v1 (CC BY 4.0)}
\begin{document}
\maketitle

\noindent\emph{Source and licence.} The Damped Waves Equation and generalized Cosine and Sine families on Banach spaces. Version arXiv:2607.05856v1 (CC BY 4.0), distributed by arXiv under the Creative Commons Attribution 4.0 International licence, http://creativecommons.org/licenses/by/4.0/, as recorded in the arXiv licence metadata for this exact version and shown on the abstract page https://arxiv.org/abs/2607.05856v1. Full licence notice: \texttt{licenses/arxiv-2607.05856-CC-BY-4.0.txt}.\par\smallskip

\noindent\textit{Editorial scope.} Counted unit: the article's lemma l2.1 (source0.tex lines 741--762). The article's complete original proof of it is delivered with it (source0.tex lines 763--875, one \texttt{proof} environment of 113 lines). No mathematics is written, repaired or re-derived: the statement and the proof are the article's own lines, in the article's own notation, and no retained line is substituted or re-flowed.  No further source-local context is needed: every cross-reference of every retained line resolves inside the two delivered spans.  Nothing was omitted from any retained span, so the package carries no omission record.  No retained line contains a citation, so the wrapper carries no bibliography and no bibliography entry.  No retained line contains a figure environment, an image-inclusion command or drawing code, so no image is delivered and the package declares no asset dependency.  Editorial formatting only, in addition: an AMS \texttt{amsart} preamble that restates the article's own declarations and macros used by the retained lines, namely the environments \texttt{lemma} on separate counters, together with the standard packages those lines need.  The retained environments print in this document as Lemma 1 in order of appearance, whereas the article prints the same items with the numbering of its own full document.  The complete native source of the article as distributed in the arXiv e-print for this exact version is bundled unchanged as \texttt{sources/204578/source0.tex}.
\begin{lemma}\label{l2.1}
	Assume $A_0:\dom(A_0)\subseteq\bX\to\bX$ and $B_0:\dom(B_0)\subseteq\bX\to\bX$ are generators
	of $C_0$-groups $\{e^{tA_0}\}_{t\in\R}$ and $\{e^{tB_0}\}_{t\in\R}$ on $\bX$, respectively, such that
	\begin{enumerate}
		\item[(i)] $A_0 B_0 = B_0 A_0$,
		\item[(ii)] $\dom(A_0)\subseteq\dom(B_0)$, $B_0(\dom(A_0^2))\subseteq\dom(A_0)$,
		\item[(iii)] $A_0^\pm := \pm A_0 - B_0$ generate $C_0$-groups $\{e^{tA_0^\pm}\}_{t\in\R}$ on $\bX$, respectively.
	\end{enumerate}
	Then, $G_0:\dom(A_0^2)\times\dom(A_0)\subseteq\dom(A_0)\times \bX\to\dom(A_0)\times \bX$ defined by
	\begin{equation}\label{eq:3}
		G_0 = \begin{bmatrix} -B_0 & I \\ A_0^2 & -B_0 \end{bmatrix}
	\end{equation}
	is the generator of the $C_0$-group $\{\cT_0(t)\}_{t\in\R}$ on $\dom(A_0)\times\bX$ given by
\begin{equation}\label{eq:4}
\cT_0(t) = \begin{bmatrix}
\tfrac{1}{2}(e^{tA_0^+}+e^{tA_0^-})\big|_{\dom(A_0)} &
\displaystyle\int_0^t e^{(t-\tau)A_0^+}e^{\tau A_0^-}\,\rmd\tau \\
\tfrac{1}{2}(e^{tA_0^+}-e^{tA_0^-})A_0\big|_{\dom(A_0)} &
\tfrac{1}{2}(e^{tA_0^+}+e^{tA_0^-})
\end{bmatrix},\; t\in\RR.
\end{equation}
\end{lemma} 

\begin{proof} First, we prove that $G_0$ and $\cT_0(t)$ are well-defined. Indeed, assumption (ii) implies that $G_0$ is well-defined and we may set $\dom(A_0^\pm)=\dom(A_0)$. Since $e^{tA_0^\pm}(\dom(A_0^\pm))\subseteq\dom(A_0^\pm)$, we have
\begin{equation}\label{l2.1.1}
		e^{tA_0^\pm}(\dom(A_0))\subseteq\dom(A_0)\quad\text{for any }t\in\R.
	\end{equation}
From assumption (i) and the definition of $A_0^\pm$ we infer that $A_0,B_0,A_0^+,A_0^-$ commute. Similarly, $ R(\lambda_1,A_0)$, $R(\lambda_2,B_0)$, $R(\lambda_3,A_0^+)$ and $R(\lambda_4,A_0^-)$ commute 
for any $\lambda_1\in\rho(A_0)$, $\lambda_2\in\rho(B_0)$, $\lambda_3\in\rho(A_0^+)$ and $\lambda_4\in\rho(A_0^-)$. It follows that 
	\begin{equation}\label{l2.1.2}
		e^{t_1A_0},e^{t_2B_0},e^{t_3A_0^+},e^{t_4A_0^-}\;\mbox{commute for any}\;t_1,t_2,t_3,t_4\in\R,
	\end{equation}	
	\begin{equation}\label{l2.1.3}
		e^{tA_0}e^{sB_0}=e^{sB_0}e^{tA_0},\;\;e^{tA_0^\pm}=e^{tA_0}e^{-tB_0}=e^{-tB_0}e^{tA_0}\;\mbox{for any}\; t,s\in\R
	\end{equation}	
	From \eqref{l2.1.1}, \eqref{l2.1.2} and \eqref{l2.1.3} we obtain 
	\begin{equation}\label{l2.1.4}
		\int_0^t e^{(t-\tau)A_0^+}e^{\tau A_0^-}x\,\rmd\tau
		= e^{tA_0^+}\int_0^t e^{-A_0\tau}e^{\tau B_0}e^{-A_0\tau}e^{-\tau B_0}x\,\rmd\tau=e^{tA_0^+}\int_0^t e^{-2\tau A_0}x\,\rmd\tau\in\dom(A_0)
	\end{equation}
	for any $t\in\R$, $x\in\bX$. Hence $\cT_0(t)$ is well-defined.
	
	Next, we prove that $\{\cT_0(t)\}$ is a $C_0$-group on $\dom(A_0)\times \bX$. Since $A_0$ and $A_0^\pm$ commute,
	\begin{equation}\label{l2.1.5}
		A_0 e^{tA_0^\pm}x = e^{tA_0^\pm}A_0 x\;\mbox{for any}\;\,t\in\R,\;x\in\dom(A_0).
	\end{equation}
	Moreover, from \eqref{l2.1.3}, \eqref{l2.1.4} and \eqref{l2.1.5} we can see that 
	\begin{equation}\label{l2.1.6}
		A_0\int_0^t e^{(t-\tau)A_0^+}e^{\tau A_0^-}x\,\rmd\tau
		= e^{tA_0^+}A_0\int_0^t e^{-2\tau A_0}x\,\rmd\tau
		= \tfrac{1}{2}(e^{tA_0^+}-e^{tA_0^-})x
	\end{equation}
	for any $x\in \bX$, $t\in\R$. Clearly, by letting $t=0$ in \eqref{eq:4}, we have $\cT_0(0) = I_{\dom(A_0)\times \bX}$. The group property of $\cT_0$ follows after a long but straightforward computation based on \eqref{l2.1.2}--\eqref{l2.1.6}. 
	
	Next, we prove the strong continuity of $\cT_0$. Fix $y\in\dom(A_0)$, $x\in \bX$. Then, from \eqref{eq:4} and \eqref{l2.1.6} we obtain 
	\begin{align}\label{2.1.14}
		&\left\|\cT_0(t)\tbinom{y}{x}-\tbinom{y}{x}\right\|_{\dom(A_0)\times\bX}=\left\|\tfrac{1}{2}(e^{tA_0^+}+e^{tA_0^-})y-y
		+e^{tA_0^+}\int_0^t e^{-2\tau A_0}x\,d\tau\right\|+\left\|\tfrac{1}{2}(e^{tA_0^+}-e^{tA_0^-})A_0 y\right\|
		\nonumber\\
		&\qquad\qquad\qquad\qquad+\left\|\tfrac{1}{2}(e^{tA_0^+}+e^{tA_0^-})A_0y-A_0y
		+\tfrac{1}{2}(e^{tA_0^+}-e^{tA_0^-})x\right\|   
		+\left\|\tfrac{1}{2}(e^{tA_0^+}+e^{tA_0^-})x-x\right\|
	\end{align}
	for any $t\in\RR$. Since $\{e^{tA_0^\pm}\}_{t\in\RR}$ and $\{e^{tA_0}\}_{t\in\RR}$ are $C_0$-groups we know that
	\begin{equation}\label{2.1.15}
		\lim_{t\to0}e^{tA_0^\pm}z_1=z_1,\quad\lim_{t\to0}\int_0^t e^{-2\tau A_0}z_2\,d\tau=0
		\;\mbox{for any}\;z_1,z_2\in\bX,
	\end{equation}
	and so $\{\cT_0(t)\}_{t\in\RR}$ is a $C_0$-group.
	
	To finish the proof of the lemma we prove that $G_0$, defined in \eqref{eq:3}, is the generator of $\{\cT_0(t)\}_{t\in\R}$. First, we show that there exists $\omega_0\in\RR$ such that $\CC_{\omega_0}^+\subseteq\rho(G_0)$ and compute $R(\lambda,G_0)$ for $\lambda\in\CC_{\omega_0}^+$. Since $\{e^{tA_0^\pm}\}$ are $C_0$-groups, there exist $\omega_0^\pm\in\R$ and $M_0^\pm\geq 0$ such that
	\begin{equation}\label{2.1.16}
		\|e^{tA_0^\pm}\|\le M_0^\pm e^{\omega_0^\pm|t|}\;\;\mbox{for any}\;\;t\in\R.
	\end{equation}
	Define $\omega_0=\max\{\omega_0^+,\omega_0^-\}$ and $Q_0(\lambda):\dom(A_0^2)\subseteq\bX\to \bX$ by
	\begin{equation}\label{2.1.17}
		Q_0(\lambda)=(\lambda I+B_0)^2-A_0^2,\;\mbox{for}\;\lambda\in\CC.
	\end{equation}
	From assumption (ii),  $\dom(A_0^2)\subseteq\dom(B_0^2)$. By assumption (i), 
	\begin{equation}\label{2.1.18}
		Q_0(\lambda)=(\lambda I+B_0-A_0)(\lambda I+B_0+A_0)=(\lambda I-A_0^+)(\lambda I-A_0^-)\;\mbox{for any}\;\lambda\in\CC.
	\end{equation}
	It follows that $Q_0(\lambda)$ is invertible with bounded inverse for $\lambda\in\mathbb{C}_{\omega_0}^+$ and 
	\begin{equation}\label{2.1.19}
		Q_0^{-1}(\lambda)=R(\lambda,A_0^-)R(\lambda,A_0^+)=R(\lambda,A_0^+)R(\lambda,A_0^-)\;\mbox{for any}\;\lambda\in\CC.
	\end{equation}
	Next, we solve 
	\begin{equation}\label{2.1.20}
		\lambda\binom{u}{v}-G_0\binom{u}{v}=\binom{f}{g},\;\mbox{with}\; u\in\dom(A_0^2)\;\mbox{and}\; v,f\in\dom(A_0), g\in\bX,\;\mbox{for}\; \lambda\in\mathbb{C}_{\omega_0}^+. 
	\end{equation}
	By \eqref{eq:3} equation \eqref{2.1.20} is equivalent to
	\begin{equation}\label{2.1.21}
		\begin{cases}v=(\lambda I+B_0)u-f\\Q_0(\lambda)u=(\lambda I+B_0)f+g\end{cases}
	\end{equation}
	By assumption (ii), $(\lambda I+B_0)$ commutes with $Q_0^{-1}(\lambda)$. Thus, the solution of equation \eqref{2.1.20} is given by
	\begin{equation}\label{2.1.22}
		\begin{cases}
			u=(\lambda I+B_0)Q_0^{-1}(\lambda)f+Q_0^{-1}(\lambda)g,\\
			v=(\lambda I+B_0)^2Q_0^{-1}(\lambda)f-f+(\lambda I+B_0)Q_0^{-1}(\lambda)g.
		\end{cases}
	\end{equation}
	We conclude that $\lambda I_{\dom(A_0)\times \bX}-G_0$ is invertible for any $\lambda\in\mathbb{C}_{\omega_0}^+$. Next, we will show that $\big(\lambda I_{\dom(A_0)\times \bX}-G_0\big)^{-1}$ is bounded on $\dom(A_0)\times\bX$. First, using \eqref{2.1.19},
	\begin{align}\label{2.1.23}
		(\lambda I+B_0)Q_0^{-1}(\lambda)&=\tfrac{1}{2}\Big((\lambda I-A_0^+)+(\lambda I-A_0^-)\Big)R(\lambda,A_0^+)R(\lambda,A_0^-)\nonumber\\
		&=\tfrac{1}{2}\Big(R(\lambda,A_0^+)+R(\lambda,A_0^-)\Big)\;\mbox{for any}\;\lambda\in\CC_{\omega_0}^+.
	\end{align}	
	Since $(\lambda I-A_0^\pm)R(\lambda,A_0^\pm)=I$ for any $\lambda\in\CC_{\omega_0}^+$, 
	\begin{equation}\label{2.1.24}
		(\lambda I+B_0)R(\lambda,A_0^+)=I+A_0 R(\lambda,A_0^+),\;(\lambda I+B_0)R(\lambda,A_0^-)=I-A_0 R(\lambda,A_0^-)
	\end{equation}
	for any $\lambda\in\CC_{\omega_0}^+$, which allows us to conclude that 
	\begin{align}\label{2.1.25}
		(\lambda I+B_0)^2Q_0^{-1}(\lambda)-I&=\tfrac{1}{2}\Big((\lambda I+B_0)R(\lambda,A_0^+)+(\lambda I+B_0)R(\lambda,A_0^-)\Big)-I\nonumber\\
		&= \tfrac{1}{2}A_0\Big( R(\lambda,A_0^+)-R(\lambda,A_0^-)\Big)\quad\text{for any }\lambda\in\mathbb{C}_{\omega_0}^+.
	\end{align}
	From \eqref{2.1.22} and \eqref{2.1.25} we obtain 
	\begin{equation}\label{2.1.26}
		(\lambda I_{\dom(A_0)\times\bX}-G_0)^{-1}
		=\begin{bmatrix}
			\tfrac{1}{2}\Big(R(\lambda,A_0^+)+R(\lambda,A_0^-)\Big)\big|_{\dom(A_0)} &
			R(\lambda,A_0^+)R(\lambda,A_0^-)\\
			\tfrac{1}{2}\big(R(\lambda,A_0^+)-R(\lambda,A_0^-)\big)A_0\big|_{\dom(A_0)} &
			\tfrac{1}{2}\big(R(\lambda,A_0^+)+R(\lambda,A_0^-)\big)
		\end{bmatrix}
	\end{equation}
	for any $\lambda\in\CC_{\omega_0}^+$. Since $\dom(A_0^\pm)=\dom(A_0)$ and $A_0$ is closed, we infer that
	$A_0 R(\lambda,A_0^\pm)\in\mathcal{B}(\bX)$ and $\Big(R(\lambda,A_0^+)+R(\lambda,A_0^-)\Big)\big|_{\dom(A_0)}\in\mathcal{B}(\dom(A_0))$ for any $\lambda\in\CC_{\omega_0}^+$. It follows that $A_0R(\lambda,A_0^+)R(\lambda,A_0^-)\in\mathcal{B}(\bX)$, hence $R(\lambda,A_0^+)R(\lambda,A_0^-)\in\mathcal{B}(\bX,\dom(A_0))$. Using the fact that $A_0\big|_{\dom(A_0)}\in\mathcal{B}(\dom(A_0),\bX)$ and $R(\lambda,A_0^\pm)\in\mathcal{B}(\bX)$ we have $\big(R(\lambda,A_0^+)-R(\lambda,A_0^-)\big)A_0\big|_{\dom(A_0)}\in\mathcal{B}(\dom(A_0),\bX)$. We conclude that all four blocks in \eqref{2.1.26} are bounded, which proves 
	$(\lambda I_{\dom(A_0)\times \bX}-G_0)^{-1}\in\mathcal{B}(\dom(A_0)\times \bX)$ for any $\lambda\in\CC_{\omega_0}^+$, thus $\lambda\in\CC_{\omega_0}^+\subset\rho(G_0)$.
	
	Taking Laplace transform in \eqref{eq:4} and using \eqref{2.1.26}, we conclude that
	\begin{equation}\label{2.1.27}
		\mathcal{L}\Bigl\{\cT_0(\cdot)\tbinom{y}{x}\Bigr\}(\lambda)=R(\lambda,G_0)\tbinom{y}{x}
		\;\mbox{for any}\;\lambda\in\mathbb{C}_{\omega_0}^+,\;y\in\dom(A_0),\;x\in\bX,
	\end{equation}
	which proves that $G_0$ is the generator of $\{\cT_0(t)\}_{t\in\R}$.
\end{proof}

\end{document}
